\chapter{第10章：模性定理与费马大定理 习题}
\difficultykey

\section{基础题 \easy}

\begin{problem}{指数四的化归}
\easy
说明若存在整数解
\[
a^n+b^n=c^n,
\qquad n>2,
\]
则只需分别处理 $n=4$ 与奇素数指数即可。特别地，说明若 $x^4+y^4=z^4$ 有非零整数解，则方程
\[
X^4+Y^4=Z^2
\]
也有非零整数解。
\end{problem}
\begin{solution}
若 $n$ 有奇素因子 $p$，写 $n=pm$，则
\[
(a^m)^p+(b^m)^p=(c^m)^p,
\]
于是可把问题化到奇素数指数 $p$。若 $n$ 是 $2$ 的幂且 $n>2$，则 $4\mid n$，写 $n=4m$，便有
\[
(a^m)^4+(b^m)^4=(c^m)^4.
\]
因此只需处理 $n=4$ 与奇素数指数。

若 $x^4+y^4=z^4$ 有解，则令
\[
X=x,\quad Y=y,\quad Z=z^2,
\]
便得到
\[
X^4+Y^4=Z^2.
\]
所以证明 $X^4+Y^4=Z^2$ 无非零整数解即可推出 $n=4$ 的费马大定理。
\end{solution}

\begin{problem}{原始解的奇偶性}
\easy
设
\[
X^4+Y^4=Z^2
\]
有一组非零整数解。证明可取一组原始解满足 $\gcd(X,Y)=1$，且 $X,Y$ 中恰有一个是偶数。
\end{problem}
\begin{solution}
若 $d=\gcd(X,Y)>1$，则 $d^4\mid Z^2$，从而 $d^2\mid Z$。把三者同除以适当公因子，可得到原始解。

在原始解中，$X,Y$ 不可能同为偶数；若同为奇数，则
\[
X^4\equiv Y^4\equiv1\pmod{16},
\]
于是
\[
Z^2=X^4+Y^4\equiv2\pmod{16},
\]
但平方数模 $16$ 不可能同余于 $2$，矛盾。故恰有一个是偶数，一个是奇数。
\end{solution}

\begin{problem}{勾股参数化的应用}
\easy
设 $(X,Y,Z)$ 是上一题中的原始解，并设 $Y$ 为偶数。说明三元组
\[
(X^2,Y^2,Z)
\]
是原始勾股三元组，因此存在互素整数 $u>v>0$，一奇一偶，使得
\[
X^2=u^2-v^2,
\qquad
Y^2=2uv,
\qquad
Z=u^2+v^2.
\]
\end{problem}
\begin{solution}
由
\[
X^4+Y^4=Z^2
\]
可知
\[
(X^2)^2+(Y^2)^2=Z^2.
\]
因为 $\gcd(X,Y)=1$，故
\[
\gcd(X^2,Y^2)=1,
\]
于是 $(X^2,Y^2,Z)$ 是原始勾股三元组。对原始勾股三元组的标准参数化给出
\[
X^2=u^2-v^2,
\qquad
Y^2=2uv,
\qquad
Z=u^2+v^2,
\]
其中 $u,v$ 互素，且一奇一偶。
\end{solution}

\begin{problem}{无穷递降的完成}
\easy
沿用上题记号，完成 Fermat 对 $n=4$ 的无穷递降证明：从一组最小的 $Z$ 出发构造出更小的正整数解，导出矛盾。
\end{problem}
\begin{solution}
取使 $Z$ 最小的原始解，并设 $Y$ 为偶数。由上题
\[
Y^2=2uv,
\qquad \gcd(u,v)=1.
\]
因为 $2uv$ 是平方数且 $u,v$ 互素，所以必可写成
\[
u=m^2,
\qquad v=2n^2
\]
或交换次序；改记后可设如此。于是
\[
X^2=u^2-v^2=m^4-4n^4=(m^2-2n^2)(m^2+2n^2).
\]
这两个因子互素且都是奇数，它们的乘积是平方，因此各自都是平方：
\[
m^2-2n^2=r^2,
\qquad m^2+2n^2=s^2.
\]
于是
\[
s^2-r^2=4n^2,
\qquad
\frac{s-r}{2}\cdot \frac{s+r}{2}=n^2.
\]
两因子互素，所以各自为平方。设
\[
\frac{s-r}{2}=a^2,
\qquad
\frac{s+r}{2}=b^2,
\qquad b>a>0.
\]
则
\[
r=b^2-a^2,
\qquad s=a^2+b^2.
\]
再由
\[
m^2=\frac{r^2+s^2}{2}
\]
得到
\[
m^2=a^4+b^4.
\]
这说明 $(a,b,m)$ 又给出一组方程
\[
a^4+b^4=m^2
\]
的正整数解，而且
\[
m< m^2 < u^2+v^2=Z.
\]
于是得到比最小 $Z$ 更小的解，矛盾。故不存在非零整数解，这就证明了 $n=4$ 的费马大定理。
\end{solution}

\begin{problem}{Frey曲线的定义与判别式}
\easy
对奇素数 $p\ge5$，设存在原始解
\[
a^p+b^p=c^p.
\]
定义 Frey 曲线
\[
E_{a,b,c}:y^2=x(x-a^p)(x+b^p).
\]
计算它的判别式，并说明为何它非奇异。
\end{problem}
\begin{solution}
把三次多项式的三根写成
\[
0,\quad a^p,\quad -b^p.
\]
若椭圆曲线写成 $y^2=(x-e_1)(x-e_2)(x-e_3)$，则判别式为
\[
\Delta=16\prod_{i<j}(e_i-e_j)^2.
\]
因此
\[
\Delta(E_{a,b,c})=16\,a^{2p}b^{2p}(a^p+b^p)^{2}=16(a b c)^{2p}.
\]
因为假设 $abc\neq0$，所以 $\Delta\neq0$，曲线非奇异。
\end{solution}

\begin{problem}{奇素数处的乘法约化}
\easy
设 $q$ 是整除 $abc$ 的奇素数。说明 Frey 曲线 $E_{a,b,c}$ 在 $q$ 处是坏约化，并解释为何这里出现的是乘法约化而不是加性约化。
\end{problem}
\begin{solution}
由上题判别式公式，若 $q\mid abc$，则
\[
q\mid \Delta(E_{a,b,c}),
\]
所以在 $q$ 处是坏约化。

设 $A=a^p$，$B=b^p$。对模型
\[
y^2=x(x-A)(x+B)
\]
有
\[
c_4=16(A^2+AB+B^2).
\]
若 $q\mid a$ 而 $q\nmid b$，则模 $q$ 下 $A\equiv0$，$B\not\equiv0$，从而
\[
c_4\equiv16B^2\not\equiv0\pmod q.
\]
这说明 $q\mid\Delta$ 但 $q\nmid c_4$，正是乘法约化的判别标准。$q\mid b$ 或 $q\mid c$ 时同理。
\end{solution}

\begin{problem}{半稳定性的陈述}
\easy
说明 Frey 曲线在费马方程应用中为何被称为半稳定曲线。
\end{problem}
\begin{solution}
前一题表明，所有整除 $abc$ 的奇素数处都只有乘法约化。对素数 $2$ 的分析较细，但结论仍是：Frey 曲线经过适当最小化后在 $2$ 处也不会出现真正的加性坏约化。

因此它在每个坏素数处不是好约化就是乘法约化，这正是“半稳定”的定义。费马方程与 Wiles 论证之所以配合得好，一个关键原因就是 Frey 曲线落在半稳定范围内。
\end{solution}

\begin{problem}{模性与水平下降的联锁}
\easy
设 Frey 曲线 $E_{a,b,c}$ 是模的，且某个模 $\ell$ 残余表示不可约。说明把“模性”与“Ribet 水平下降”连起来后会得到怎样的总体结论。
\end{problem}
\begin{solution}
模性说明 $\bar\rho_{E_{a,b,c},\ell}$ 来自某个权 $2$ 新形式，其级别等于 Frey 曲线的导子。Ribet 水平下降进一步说明：若残余表示不可约，则可以把级别中由半稳定坏约化带来的素因子逐个去掉。

对费马方程对应的 Frey 曲线，这一过程会把级别降到极小的值，最终降到 $2$。于是若 level $2$ 上不存在相应新形式，就得到矛盾。
\end{solution}

\begin{problem}{级别二空间的消失}
\easy
说明为什么一旦把问题降到权 $2$、级别 $2$ 的新形式，就会立刻结束证明。
\end{problem}
\begin{solution}
关键事实是
\[
S_2(\Gamma_0(2))=0.
\]
也就是说，级别 $2$ 的权 $2$ 尖点形式根本不存在，更不用说归一化新形式了。

因此如果某个残余表示被证明必须来自 level $2$ 的权 $2$ 新形式，就立刻产生矛盾。这正是 Frey 曲线、模性提升与 Ribet 水平下降共同指向的终点。
\end{solution}

\section{进阶题 \medium}

\begin{problem}{Langlands和Tunnell起点}
\medium
说明在 Wiles 证明中，Langlands--Tunnell 定理扮演什么角色。
\end{problem}
\begin{solution}
Wiles 的模性提升并不是凭空开始的，它需要一个“起点”：某个残余表示已经知道是模的。对模 $3$ 的奇不可约表示，Langlands--Tunnell 定理给出这一起点，说明这类表示来自权 $1$ 模形式，进而可转化为适合模性提升使用的输入。

因此，Langlands--Tunnell 的作用是提供残余模性。之后再通过模性提升定理，把“残余模”升到“整个 $\ell$-进表示模”。
\end{solution}

\begin{problem}{三五开关技巧}
\medium
陈述并解释三五开关技巧的核心思想。
\end{problem}
\begin{solution}
三五开关技巧用于处理某条椭圆曲线的模 $3$ 表示可能不可约性不足的问题。其思路是：若模 $3$ 信息不好，就利用曲线的模 $5$ 挠结构构造另一条辅助椭圆曲线，使新曲线在模 $5$ 上与原曲线相关，而在模 $3$ 上具有更适合应用 Langlands--Tunnell 的性质。

这样就能在 $3$ 与 $5$ 之间切换：先从较好的那一侧得到残余模性，再借助模性提升把信息带回原曲线。它是 Wiles 论证中一个非常精巧的桥梁步骤。
\end{solution}

\begin{problem}{残余模性到提升定理}
\medium
设某条半稳定椭圆曲线 $E/\Q$ 的残余表示 $\bar\rho_{E,3}$ 已知模。说明模性提升定理要解决的目标是什么。
\end{problem}
\begin{solution}
残余模性只告诉我们
\[
\bar\rho_{E,3}
\]
来自某个模形式；但真正需要的是整个 $3$-进表示
\[
\rho_{E,3}:G_\Q\to\GL_2(\Z_3)
\]
来自模形式。模性提升定理的目标正是：在适当局部条件、不可约性和极小性假设下，从 $\bar\rho$ 的模性推出 $\rho$ 的模性。

换言之，它把“模 $3$ 的信息”提升为“全部 $3$-进信息”，从而把椭圆曲线放入模性定理的范围。
\end{solution}

\begin{problem}{变形问题的内容}
\medium
给定残余表示
\[
\bar\rho:G_\Q\to\GL_2(k),
\]
说明在 Wiles 理论中，一个“变形问题”通常要求固定哪些数据。
\end{problem}
\begin{solution}
一个变形问题通常考虑把 $\bar\rho$ 提升到某个完备局部环 $A$ 上的表示
\[
\rho_A:G_\Q\to\GL_2(A),
\]
并要求模极大理想约化后回到 $\bar\rho$。除此之外，还要固定一些局部条件，例如：在某些素数处要求未分歧、有限平坦、半稳定或极小；通常还固定行列式，使之等于分圆特征乘上已知角色。

因此，变形问题并不是“所有可能提升”的总和，而是“满足指定局部条件与行列式的提升”的集合。
\end{solution}

\begin{problem}{普遍变形环的意义}
\medium
解释什么是普遍变形环 $R$，以及它为什么能把一个表示论问题变成交换代数问题。
\end{problem}
\begin{solution}
若某个变形函子可表出，就存在一个完备 Noether 局部环 $R$ 和一个“普遍提升”
\[
\rho^{\mathrm{univ}}:G_\Q\to\GL_2(R),
\]
使任何满足条件的变形都唯一来自某个局部同态
\[
R\to A.
\]
这就是普遍变形环。

它把“分类所有提升”转化成“研究一个具体局部环的结构”。于是表示论与数论问题被翻译为关于生成元、关系式、维数与交叉数的交换代数问题，这正是 Wiles 方法的核心。
\end{solution}

\begin{problem}{Hecke代数的意义}
\medium
说明 Hecke 代数 $T$ 在 Wiles 方法中代表什么对象。
\end{problem}
\begin{solution}
Hecke 代数 $T$ 由相关模形式空间上的 Hecke 算子生成，它编码了模形式的全部 Hecke 特征值数据。对一个给定的残余表示，局部化后的 Hecke 代数对应“与该残余表示同余的模形式家族”。

因此，$T$ 可以看成“模表示一侧”的坐标环：它记录所有实际来自模形式的 Galois 表示变形。若能证明 $R=T$，就说明每个允许的变形都确实来自模形式。
\end{solution}

\begin{problem}{R等于T的含义}
\medium
解释等式
\[
R=T
\]
在模性提升中的数学含义。
\end{problem}
\begin{solution}
$R$ 参数化所有满足局部条件的 Galois 表示变形，$T$ 参数化实际由模形式产生的变形。若能构造满射
\[
R\twoheadrightarrow T
\]
并进一步证明它是同构，就说明“允许出现的变形”和“真正模的变形”完全相同。

于是任一满足该变形问题条件的 $\ell$-进表示，只要其残余表示是模的，就必然自身也是模的。这就是模性提升定理的本质内容。
\end{solution}

\begin{problem}{切空间与Selmer群}
\medium
说明为什么普遍变形环的切空间维数会与某个 Selmer 群维数相关。
\end{problem}
\begin{solution}
一阶无穷小变形对应于把系数环从 $k$ 提升到双数环
\[
k[\varepsilon]/(\varepsilon^2).
\]
这样的变形类由 Galois 上同调控制，具体地说，满足局部条件的一阶变形组成某个 Selmer 群。

因此，普遍变形环的切空间维数
\[
\dim_k \mathfrak m_R/(\mathfrak m_R^2,\varpi)
\]
可以用 Selmer 群维数表达。Wiles 方法正是通过比较 Selmer 与对偶 Selmer，来估计 $R$ 的大小与关系数。
\end{solution}

\begin{problem}{Wiles数值判据}
\medium
简述 Wiles 数值判据想要比较的两个量是什么，以及它的结论为何强大。
\end{problem}
\begin{solution}
Wiles 数值判据比较的是：一边来自 Hecke 模的代数长度或交叉数，另一边来自变形环切空间与对偶切空间的数据。粗略地说，它把“$R$ 有多少关系”和“$T$ 有多大”放到同一个数值框架中。

当这些数值恰好相等时，就能从已有的满射
\[
R\twoheadrightarrow T
\]
推出它其实是同构。这样，一个本来看似很深的“所有变形都模”命题，被压缩为可计算的数值比较。
\end{solution}

\begin{problem}{指数大于等于五的逻辑链}
\medium
把 $p\ge5$ 时费马大定理证明的逻辑链按顺序写出，不要求技术细节。
\end{problem}
\begin{solution}
逻辑链如下：
\begin{itemize}
\item[(i)] 反设存在原始解 $a^p+b^p=c^p$；
\item[(ii)] 构造 Frey 曲线 $E_{a,b,c}$，它是半稳定椭圆曲线；
\item[(iii)] 由 Wiles--Taylor--Wiles，半稳定椭圆曲线是模的；
\item[(iv)] 选择适当的残余表示并验证其不可约性；
\item[(v)] 由 Ribet 水平下降，把相应模表示的级别降到 $2$；
\item[(vi)] 但 $S_2(\Gamma_0(2))=0$，level $2$ 根本没有相应新形式；
\item[(vii)] 得到矛盾，故原始解不存在。
\end{itemize}
这就证明了所有奇素数指数的情形。
\end{solution}

\section{挑战题 \hard}

\begin{problem}{Taylor和Wiles辅助素数}
\hard
说明 Taylor--Wiles 方法中为什么要引入一批辅助素数。
\end{problem}
\begin{solution}
只研究原始变形问题时，普遍变形环的关系数往往难以控制。Taylor--Wiles 引入一批精心挑选的辅助素数 $Q$，在这些素数处放宽局部条件，从而得到一族更容易计算的扩展变形环 $R_Q$ 与相应模 $M_Q$。

这些辅助素数的选择使得对偶 Selmer 群被逐步“杀掉”，环与模的维数控制变得整齐，最终可以在极限中恢复原问题。这就是辅助素数的核心作用。
\end{solution}

\begin{problem}{Patching的核心想法}
\hard
用宏观语言解释 Taylor--Wiles patching 的基本步骤。
\end{problem}
\begin{solution}
Patching 的思路是：不直接比较原始的 $R$ 与 $T$，而是先构造一系列带辅助素数的对象
\[
R_Q,
\qquad M_Q,
\qquad T_Q.
\]
这些对象在有限层面上满足更好的生成元与关系数控制。

随后把所有层面同时“拼接”成一个极限对象，使其既保留足够多的自由性，又反映原问题的 Hecke 模结构。最后在这个 patched 世界中证明自由性，再下推回原始层面，得到 $R=T$。因此 patching 的本质是“先在更大空间里变简单，再投影回原问题”。
\end{solution}

\begin{problem}{极小与非极小变形}
\hard
解释“极小变形”和“非极小变形”的区别，以及为什么 Wiles 最先处理的是极小情形。
\end{problem}
\begin{solution}
极小变形要求在每个素数处的局部 ramification 尽量不比残余表示更大；也就是说，提升过程中不额外增加导子。非极小变形则允许在某些素数处出现更多 ramification。

极小情形的局部条件最紧，因此对应的 Selmer 群和对偶 Selmer 群最容易控制，数值判据也最干净。Wiles 先在极小情形证明 $R=T$，再通过更精细的论证逐步扩展到非极小情形。
\end{solution}

\begin{problem}{对偶Selmer群的作用}
\hard
说明对偶 Selmer 群在 Wiles 证明中为何如此关键。
\end{problem}
\begin{solution}
Selmer 群测量允许的一阶变形有多少，而对偶 Selmer 群测量这些变形之间还缺多少关系。若对偶 Selmer 群很大，普遍变形环就可能有很多关系，难以与 Hecke 代数匹配。

Taylor--Wiles 方法的关键正是利用辅助素数把对偶 Selmer 群压缩到零或足够小。这样一来，变形环的关系数得到控制，最终能证明它是完全交并与 Hecke 一侧大小一致。
\end{solution}

\begin{problem}{完全交与自由性}
\hard
解释为什么在 $R=T$ 论证中，“完全交”和“Hecke 模自由”是两个重要目标。
\end{problem}
\begin{solution}
若变形环是完全交，则它的生成元数与关系数之间达到最理想的平衡，这会极大简化维数与深度计算。另一方面，若相关 Hecke 模在该环上是自由的，就说明模形式一侧没有隐藏的扭结，结构与环本身完全匹配。

在 patched 层面证明这两点后，就能把交换代数信息精确传回原始层面，从而推出 Hecke 代数与变形环同构。简言之，完全交控制环，自由性控制模，两者合起来给出 $R=T$。
\end{solution}

\begin{problem}{两个方向的映射含义}
\hard
在很多概述中会看到
\[
R\twoheadrightarrow T
\qquad\text{和}\qquad
T\hookrightarrow R
\]
这样的说法。解释这两个方向分别表达了什么思想。
\end{problem}
\begin{solution}
满射
\[
R\twoheadrightarrow T
\]
表示：每个 Hecke 特征值系统都对应一个满足局部条件的 Galois 表示变形，因此模形式一侧嵌入到所有允许变形之中。这通常来自构造模形式关联的 Galois 表示。

而把 $T$ 看成注入某个关于 $R$ 的对象，则体现相反方向：模形式给出的变形已经“足够多”，足以探测 $R$ 的全部结构。数值判据和 patching 说明这两个方向没有缺口，最终只能是同构。
\end{solution}

\begin{problem}{半稳定模性为何足够}
\hard
解释为什么只要证明“所有半稳定椭圆曲线都是模的”，就已经足以推出费马大定理。
\end{problem}
\begin{solution}
费马方程若有反例，就会产生一条 Frey 曲线，而这条曲线正是半稳定的。因此假如我们已经知道所有半稳定椭圆曲线都模，那么 Frey 曲线自动模。

一旦 Frey 曲线模，再配合 Ribet 水平下降，就把其残余表示降到 level $2$，与
\[
S_2(\Gamma_0(2))=0
\]
矛盾。所以“半稳定模性”正好覆盖了费马反例所必然产生的全部曲线，因而已经足够证明费马大定理。
\end{solution}

\begin{problem}{费马大定理证明摘要}
\hard
给出费马大定理完整证明思路的摘要，要求同时包含 $n=4$ 与 $n\ge5$ 两部分。
\end{problem}
\begin{solution}
证明分两部分。

第一部分是 $n=4$。把问题化为
\[
X^4+Y^4=Z^2,
\]
然后对最小的正整数 $Z$ 做无穷递降：由原始解构造出更小的同类解，矛盾，因此不存在 $n=4$ 的反例。

第二部分是 $n\ge5$。若存在奇素数指数 $p$ 的反例，就构造 Frey 曲线 $E_{a,b,c}$。它是半稳定的；由 Wiles 与 Taylor--Wiles，半稳定椭圆曲线都模；由 Ribet，相关残余表示可以水平下降到 level $2$；但 level $2$ 没有权 $2$ 新形式，因此矛盾。综上，所有 $n>2$ 都没有非零整数解。
\end{solution}

\begin{problem}{证明结构的全景回顾}
\hard
从“反设存在反例”开始，用不超过八句话概括 Wiles 证明与 Ribet 定理如何拼接成一条完整路线。
\end{problem}
\begin{solution}
可以概括为：
\begin{itemize}
\item[(i)] 反设存在 $a^p+b^p=c^p$ 的原始反例，$p\ge5$ 为奇素数；
\item[(ii)] 由此构造 Frey 曲线，它是半稳定椭圆曲线；
\item[(iii)] Langlands--Tunnell 与三五开关给出残余模性的起点；
\item[(iv)] Wiles 的模性提升与 Taylor--Wiles patching 证明半稳定曲线模；
\item[(v)] 因而 Frey 曲线的残余表示来自某个权 $2$ 新形式；
\item[(vi)] Ribet 水平下降把级别降到 $2$；
\item[(vii)] 但 level $2$ 上没有权 $2$ 新形式；
\item[(viii)] 所以反例不存在，费马大定理成立。
\end{itemize}
\end{solution}
